\section{第二日第二讲}

\subsection{一个奇迹}

我们从 Riemann zeta 函数讲起，包括对全体整数 $\mathbb Z$ 的求和。某些书里会改用另一个变量 $y=x^2$ 来写下面的公式。

Riemann 证明了
\begin{equation}
  \int_0^\infty x^s\sum_{n\in\mathbb Z}e^{-\pi n^2x^2}\,d^\times x=2\pi^{-s/2}\Gamma(s/2)\zeta(s)=L(\chi,0)
  \label{eq:5-1}
\end{equation}
% [Tutor Note: 原文右端写作 $L(\chi,0)$。按上下文 $\chi:x\mapsto|x|^s$ 应是 $L(\chi,s)$，此处原文疑似笔误，译文照原文保留，已在 reports/ch05.md 登记。]
其中
\begin{equation}
  \chi:\ x\mapsto|x|^s,
  \label{eq:5-2}
\end{equation}
\begin{equation}
  \zeta(s)=\prod_p\frac{1}{1-p^{-s}},
  \label{eq:5-3}
\end{equation}
\begin{equation}
  \sum_{n\in\mathbb Z}e^{-\pi n^2x^2}=\Theta(x),
  \label{eq:5-4}
\end{equation}
即 Jacobi theta 函数，并记
\begin{equation}
  \zeta_\infty(s)=\pi^{-s/2}\Gamma(s/2).
  \label{eq:5-5}
\end{equation}
另外，只是作为一个记号，$d^\times x=\frac{dx}{x}$。

上式是 theta 级数的 Mellin 变换——把它理解为乘法 Fourier 变换。

注意对 theta 函数，在 Poisson 求和公式下有
\begin{equation}
  \theta(x)\ \overset{\text{PSF}}{\longleftrightarrow}\ x^{-1}\theta(x^{-1}),
  \label{eq:5-6}
\end{equation}
类似地
\begin{equation}
  \zeta_\infty(s)\longleftrightarrow\zeta(1-s).
  \label{eq:5-7}
\end{equation}
注意这是周期函数的一例，记作 $P_X$，其中 $X=\mathbb A^1$ 带 $G=\mathbb G_m$ 作用。

注意，把 L-函数写成周期，并没有普遍的理论。在代数几何里，我们有
\begin{equation}
  \text{``Schur Variety''}\ \overset{\text{周期}}{\cong}\ \text{``Schur variety''}.
  \label{eq:5-8}
\end{equation}
% [Tutor Note: 原文左右两侧引号内的写法只差 Variety/variety 的大小写，疑为原文笔误或打字噪声；译文照原文保留大小写差别。]

对阿代尔情形，(5.1) 推广为
\begin{equation}
  \int_{[G_m]}|x|^s\sum_{\gamma\in F=\mathbb A^1(F)}\Phi(\gamma x)\,d^\times x
  \label{eq:5-9}
\end{equation}
这里的积分区域是
\begin{equation}
  [G_m]=G(F)\backslash G(\mathbb A)\quad\text{或}\quad G(F)\backslash G(\mathbb A)/G(O).
  \label{eq:5-10}
\end{equation}
\begin{equation}
  P_X(x)\ \sum_{\gamma\in F=\mathbb A^1(F)}\Phi(\gamma x)\in S(X(\mathbb A))
  \label{eq:5-11}
\end{equation}
% [Tutor Note: 原文此式左端 $P_X(x)$ 与求和号之间没有等号，译文照原文保留。]
是周期，其中 $S(X(\mathbb A))$ 表示 Schwartz 函数，$X$ 是局部的，而
\begin{equation}
  \Phi(\gamma x)=\prod_{p<\infty}\mathbf 1_{\mathbb Z_p}\otimes e^{-\pi x^2_\infty},
  \label{eq:5-12}
\end{equation}
其中 $\mathbf 1_{\mathbb Z_p}$ 带 $G(\hat{\mathbb Z})$ 的单位作用。

这里我们用 idele 类特征
\begin{equation}
  F^\times\backslash\mathbb A^\times\to\mathbb C^\times.
  \label{eq:5-13}
\end{equation}
下面给几条注解。

\begin{enumerate}[label=(\arabic*),leftmargin=2.0em,itemsep=0.25em]
  \item 上述两者的不同在于：一个是自守形式，一个是自守函数。
  \item 可以把上述看作一个带两条边界的 TQFT：一条带 $G(\mathbb R)$ 作用，一条带 $G(\mathbb Q_p)$ 作用。
  \item 当 $X$ 在 $O_v$ 上光滑仿射时，$\mathbf 1_{X(O_v)}$ 是基本函数。
  \item 这个积分是下述构造的特例
  \begin{equation}
    \langle\text{自守形式},\text{周期}\rangle=\text{L-函数}
    \label{eq:5-14}
  \end{equation}
  \item 5.6 可以看作在 $T^*\mathbb A^1=M$ 的两个子空间之间作对换。也就是说，在 PSF 下，
  \begin{equation}
    \mathbb A^1\longleftrightarrow O(\mathbb A^1)^*
    \label{eq:5-15}
  \end{equation}
  在 $T^*\mathbb A^1$ 里是不变的。这里 $M$ 是一个 Hamilton $G$-空间。
\end{enumerate}

现在可以回答这个问题：Riemann 的构造为什么奏效？把 Tate 积分展开（忽略 $\gamma=0$），得
\begin{equation}
  \int_{\mathbb A^1}|x|^s\Phi(\gamma x)\,d^\times x=\prod_{p\le\infty}\int_{\mathbb Q_p^\times}\Phi_p(x)\,d^\times x.
  \label{eq:5-16}
\end{equation}
固定 $p<\infty$，我们有
\begin{equation}
  \int_{\mathbb Z=\mathbb Q_p^\times/\mathbb Z_p^\times}|x|^s\mathbf 1_{\mathbb Z_p^\times}(x)\,d^\times x=\sum_{i\ge0}p^{-is}=\frac{1}{1-p^{-s}}=\zeta_p(s).
  \label{eq:5-17}
\end{equation}
这就是一个奇迹！！！另外，这里 $i=0$ 这一项为 $1$。

\subsection{基本函数的 Mellin 变换}

设 $N:=S(\mathbb Q_p)^{\mathbb Z_p^\times}$。它是 Hecke 代数
\[
  H(\mathbb Q_p^\times,\mathbb Z_p^\times)=\mathbb C[\mathbb Q_p^\times/\mathbb Z_p^\times]=\mathbb C[\mathbb Z]
\]
的模。设 $\mathbf 1_{\mathbb Z_p}\in N$ 是 $\mathbb Z_p$ 上的常函数。

$\zeta_p(s)$ 会出现在 $N$ 的模结构里吗？答案是不会。$N$ 是由 $\mathbf 1_{\mathbb Z_p}$ 生成的秩 $1$ 自由 $\mathbb Z$-模，但考虑
\[
  N^*:=S(\mathbb Q_p^\times)^{\mathbb Z_p^\times}=\mathbb C[\mathbb Z],
\]
它包含元素 $\mathbf 1_{\mathbb Z_p^\times}$，也带 $H(\mathbb Q_p^\times,\mathbb Z_p^\times)$ 的作用，而 $\mathbf 1_{\mathbb Z_p^\times}$ 的 Mellin 变换是 $1$ 而不是 $\zeta_p(s)$。

只有把这个模看作内积空间时，$\zeta$ 才开始出现。内积 $\langle\ ,\ \rangle$ 来自嵌入 $N\hookrightarrow L^2(\mathbb Q_p)$，也就是 $L^2(\mathbb Q_p^\times)$。

澄清一个吹毛求疵的点：做积分时，我们其实考虑的是 $\mathbf 1_{\mathbb Z_p}(x)|dx|^{1/2}$，其中 $|dx|^{1/2}$ 称为半密度。

\subsection{Plancherel 公式}

我们想要形如
\begin{equation}
  \langle f_1,f_2\rangle=\langle\hat f_1,\hat f_2\rangle
  \label{eq:5-18}
\end{equation}
的公式，其中 $\hat f$ 是 $f$ 的 Mellin 变换。函数定义在 $\mathbb Q_p^\times/\mathbb Z_p^\times\equiv\mathbb Z$ 上，所以它们的 Mellin 变换定义在 $\mathbb Z$ 的对偶上；这个对偶是 $S^1\subset\mathbb C$，即酉特征标群。

可以算出
\[
  \widehat{\mathbf 1_{\mathbb Z_p}|dx|^{1/2}}(z)=\frac{1}{1-p^{-1/2}z}\,|d^\times z|^{1/2},
\]
这里 $z\in S^1$ 对应于此前的 $p^{-s}$，而 $p^{-1/2}$ 的移位源于半密度 $|dx|^{1/2}$。

此时公式 (5.18) 读作
\[
  \begin{aligned}
  &\langle\mathbf 1_{\mathbb Z_p}|dx|^{1/2},\mathbf 1_{\mathbb Z_p}|dx|^{1/2}\rangle\\
  &\qquad=\int_{S^1}\frac{1}{(1-p^{-1/2}z)(1-p^{-1/2}z^{-1})}\,d^\times z.
  \end{aligned}
\]
也可以用 Hecke 代数的作用来做。取 $h\in H=\mathbb C[\mathbb Z]$，$h$ 作用在第一项上会在右端把被积函数乘以 $\hat h$。即
\[
  \begin{aligned}
  &\langle h\cdot\mathbf 1_{\mathbb Z_p}|dx|^{1/2},\mathbf 1_{\mathbb Z_p}|dx|^{1/2}\rangle\\
  &\qquad=\int_{S^1}\frac{\hat h(z)}{(1-p^{-1/2}z)(1-p^{-1/2}z^{-1})}\,d^\times z\\
  &\qquad=\int_{S^1}\hat h(z)\zeta_p(\tfrac12+s)\zeta_p(\tfrac12-s)\,d^\times z.
  \end{aligned}
\]

这里藏着对偶的 $\mathbb G_m$，即原 $\mathbb G_m$ 的 Langlands 对偶，以及对偶 $T^*\mathbb A^1=:\check M$。等我们为 L-函数建立一般的形式化时，我们会说这里的 zeta 函数与这个对偶 $\mathbb G_m$ 在同一空间上的作用相关联——这纯属巧合。但这只是对偶性一个非常特殊的例子。

\subsection{Satake 同构}

现在回顾经典 Satake 同构。

设 $H_G=H(G(F),G(O))$ 是 $F$ 上连通分裂约化群概形 $G$ 的球 Hecke 代数。

经典 Satake 同构断言 $H_G\simeq H_T^W$，其中 $T=B/N$。由于下面这族同构，它也等价于 $H_G\simeq\mathbb C[\Rep\check G]$：
\begin{equation}
  H_G\simeq H_T^W=\mathbb C[X_\bullet(T)]^W=\mathbb C[\check T]^W=\mathbb C[\check G]^{\check G}\simeq\mathbb C[\Rep\check G].
  \label{eq:5-19}
\end{equation}
后一个同构 $\mathbb C[\check G]^{\check G}\simeq\mathbb C[\Rep\check G]$ 由特征标映射给出，即从 $\Rep\check G$ 的 Grothendieck 群到 $\check G$ 在自身上伴随作用的不变环的映射。

Hecke 同构的构造这里不回顾，我们只想指出：它用到 $H_G\otimes H_T$-双模 $S(N(F)\backslash G(F))^{G(O)}$。

也可以说，整个图景只有在几何 Satake 等价之后才完全说得通：它把 $H_G$ 变成半单范畴，并断言它等价于 $\Rep\check G$。

此外，经典 Satake 同构从分析的角度看并不很好用。为了看清楚，我们比较两侧的内积。

在右端我们有典范基 $\{s_\lambda\}$，即最低权 $\lambda$ 的不可约表示 $\{V_\lambda\}$ 的 Weyl 特征标；这里 $\lambda$ 取 $T$ 的反支配余特征标（必须取最低权 $\lambda\in X^-_\bullet(T)$，才能使下面的公式成立）。

同时，左端有一个内积，来自嵌入 $H_G\hookrightarrow L^2(G)$。

自然要问：右端的典范基是否对应左端的标准正交基？答案是否。

记 $\hat h$ 为 $h\in H_G$ 在 Satake 同构下的像，并定义 $h_\lambda$ 为满足 $\hat h_\lambda=s_\lambda$ 的元素。

左端当然有一组正交基，由同一集合 $\{\lambda\in X^-_\bullet(T)\}$ 自然地参数化：只需取双 $G(O)$-陪集的示性函数 $\mathbf 1_{G(O)\omega_\lambda G(O)}$。但这并不是我们要找的那组基：$h_\lambda\neq\mathbf 1_{G(O)\omega_\lambda G(O)}$，而且 $\{h_\lambda\}$ 关于来自 $H_G\hookrightarrow L^2(G)$ 的内积并不正交。

于是得到下面的问题：$H_G$ 上使 $h_\lambda$ 正交的自然内积是什么？

为回答它，先回忆 $s_\lambda$ 的正交性是什么意思：可以把 $\{s_\lambda\}$ 看作极大紧子群 $\check G_c\subset G$ 的特征标，于是有
\[
  \int_{\check G_c}s_\lambda\overline{s_\mu}=\delta_{\lambda\mu}.
\]
用 Weyl 积分公式可改写为
\begin{equation}
  \frac{1}{|W|}\int_{\check T_c}s_\lambda(t)\overline{s_\mu(t)}\,d_{Weyl}(t)
  =\int_{\check G_c}s_\lambda\overline{s_\mu}=\delta_{\lambda\mu},
  \label{eq:5-20}
\end{equation}
其中 $d_{Weyl}(t)=\prod_{\check\alpha\in R(\check G)}(1-e^\lambda)$ 是来自 Weyl 特征标公式的 Weyl 测度（这里 $e^\lambda:T\to\mathbb C^\times$ 是由 $\lambda\in X^\bullet(T)$ 给出的映射）。
% [Tutor Note: 原文乘积的指标是 $\check\alpha$，但因子里却写 $e^\lambda$；按 Weyl 特征标公式应是 $\prod_{\check\alpha}(1-e^{\check\alpha})$。译文照原文保留。]

我们要用这个 Weyl 测度在 $G$ 上造一个使 $h_\lambda$ 正交的内积。上面的讨论表明，需要对它作某种非平凡的修正才行。结果表明，正确的修正由群 $G$ 上的 Macdonald 公式提供：我们可以定义
\begin{equation}
  \langle h,\mathbf 1_{G(O)}\rangle=\langle h*\mathbf 1_{G(O)},\mathbf 1_{G(O)}\rangle
  =\frac{1}{|W|}\int_{T_{\check G}}\frac{\hat h(t)}{\prod_{\check\alpha\in R(\check G)}(1-q^{-\frac12}e^{\check\alpha})}\,d_{Weyl}(t).
  \label{eq:5-21}
\end{equation}
右端的修正测度就是该群的 Plancherel 测度。

修正因子可以用 $\check G$ 在其 Lie 代数上伴随表示的 $q$-特征标重写。对此的概念性解释由导出 Satake 同构提供。
